% Transcription — fonds Alexandre Grothendieck, shelfmark 8, batch 3 (pages 41-60).
% Established from the facsimile archives/batches/8/batch-03.pdf, cut from a
% copy of 8.pdf fetched through the project's own relay
% (https://grothendieck-relay.onrender.com/source/8.pdf) rather than from
% grothendieck.umontpellier.fr directly; per the skill transcribe-grothendieck.
% The batch file carries no cover sheet: PDF page k is the archivists' page
% 40+k, and their pencil numbers bottom-left agree throughout. The folder is
% eighty pages; this is the third of four batches (the others are
% transcribed separately by other passes).
%
% Pass: Opus 5.5 (claude-opus-5-5), 2026-09-26 — first pass, unchecked against the pages by a human.
%
% The folder sits in the inventory group "4-9" « Cristaux »; it has its own
% date in the catalogue, which \dating{} carries verbatim, with no group
% sentence.
%
% Pages skipped: 42, 44, 46, 48, 50, 52, 54 and 56, typed leaves « n° 382 »
% (pp. 79, 78, 77, 76, 73, 43, 42, 40 of a stencilled text on C^r manifolds:
% jets and differential operators §7.1, tensor bundles §6.8-6.9, weakening
% of structure §4.13, group varieties §4.12, transversal pairs §4.11), five
% of them upside down, with nothing in his hand. They belong to the same
% stencilled text batch 2 skipped (pp. 34-40); unrelated versos, skipped by
% the settled rule. No leaf of the proper-base-change typescript found in
% batch 1 occurs in this batch.
%
% Pages summarised rather than set line by line: 57, 58, 59, pencil graphs
% on squared paper (families of polygonal lines, abscissae 1-18, ordinates
% 1-16) with his lists of fraction comparisons; the drawings are described
% in a \note and every legible inscription is transcribed.
%
% Pages transcribed from his hand: 41, 43, 45, 47, 49, 51, 53, 55, 60.
% Fast drafting in blue-black ink; formulas carry, much of the connective
% prose does not and is left \ill{}.
%
% His pagination: none on these leaves.
%
% What the batch is about. Pp. 41-55 continue batch 2's run « Détermination
% du module de Dieudonné filtré du groupe de Barsotti-Tate » (K/Q_p, A, pi,
% k = F_q, q = p^f, M = Omega (x) W, phi_nu, varpi). P. 41: Hom_A(G_0, G'_0)
% for dim G_0 = 1 (nonzero iff pi = pi'; Hom(Omega', Omega) or
% pi Hom(Omega', Omega) according to nu, nu'); the general case up to
% isogeny, F_M^f = varpi (x) id, the « effectivity » condition (varpi and
% p^f varpi^{-1} each leave a lattice invariant). P. 43: omega_{G_0},
% the case nu = 0, Omega = A, G_0 = G_0(A, pi), M = A (x) W,
% F_M = id (x) sigma + pr_0((pi - 1) (x) sigma); the project of lifting G_0
% to a BT group G over A with A-action inducing the obvious action on t_G
% (Lubin-Tate). Pp. 45-47: via the deformation theory (divided powers on
% pi A), lifting = a filtration of Omega (x)_{Z_p} A by A (x) A-modules;
% exact sequences 0 -> omega_G -> Omega (x) A -> t_{G*} -> 0 and its
% transpose (**), t_G = quotient of Omega^v (x) A by base change along
% p_0 : A (x) A -> A, whence t_G = Omega^v, omega_G = Omega (Lubin-Tate
% existence and uniqueness, weak form); without divided powers, the
% sequences over K. P. 49: t_G = pi^{-alpha} Omega^v, D*(G)_K = M_K, and the
% question whether the two lattices D*(G), M coincide (which would force
% alpha = 0). Pp. 51-55: Tate's theorem for the Lubin-Tate group: K_pi,
% generator e of T_p(G)*, T_p (x)_{Z_p} C = prod_sigma C_sigma with Galois
% action sigma(chi_pi(g)), compared with Tate's decomposition (**)
% T_p (x) C = t (x)_A C(1) + t^v* (x)_A C, which amounts to elements x_sigma
% of C with g(x_sigma) = sigma(chi_pi(g)) x_sigma (sigma != 1),
% g(x_1) = tau(g)^{-1} chi_pi(g) x_1; there is no x'_1 with
% g(x'_1) = chi_pi(g) x'_1 (a norm argument). Pp. 57-59: pencil graphs
% (slopes n/17 ...). P. 60: the Galois group of a Tannakian category C
% with a fibre functor, example of F-crystals (filtered) up to isogeny over
% a complete DVR V of unequal characteristic with perfect residue field k;
% T_p : C -> Mod_f(pi, Q_p), C = Rep_{Q_p}(G) with G pro-algebraic (?).
%
% Notation, aligned with batches 1 and 2: \varpi for his « curly pi »
% (F^f = varpi), \Omega for the invertible A-module, \check\Omega for its
% dual, \mathbb{D}^*(G) for his D*(G) (p. 49), W, K_0 as he writes them;
% \boldsymbol{\pi} for his bold pi naming Gal(\bar K/K) (pp. 51-55), distinct
% from the uniformiser \pi; \mathfrak{t} for the gothic-looking letter of
% pp. 53-55 in Tate's decomposition (doubtful reading, noted once); \Phi
% for his barred-circle letter on p. 60 (the fibre functor).
%
% Hardest pages and readings to check first: p. 43 middle (three
% overwritten lines around « La filtration canonique »); p. 45 top and
% bottom; p. 47, the base-change subscript in the last formula; p. 51,
% « Nous supposons fixé K_pi ... »; p. 53, the letter read \mathfrak{t};
% p. 55, the line « Mais la ... »; p. 57, several fractions; p. 60, the
% annotations « gr. déri. » and the word after « pro- ».

\documentclass[11pt,a4paper]{article}
% Le préambule commun à toutes les transcriptions du fonds Grothendieck.
%
% Il tient en une idée : **l'incertitude se compose, elle ne se résout pas.**
% Une transcription qui lisse un mot douteux a détruit la seule chose pour
% laquelle on la faisait — un lecteur doit pouvoir savoir, à chaque endroit,
% ce qui était sur le feuillet et ce qui a été ajouté. Les six macros qui
% suivent sont donc l'appareil critique, et non de la décoration.
%
% Les mêmes commandes sont reconnues par `scripts/render.mjs`, qui produit la
% vue de lecture du site. Une macro ajoutée ici doit l'être aussi là-bas,
% faute de quoi le rendu échouera bruyamment — ce qui est voulu : un rendu
% silencieusement faux est pire qu'un rendu absent.

\ProvidesPackage{grothendieck}[2026/08/08 Transcriptions du fonds Grothendieck]

\usepackage{amsmath,amssymb,amsthm}
\usepackage{mathrsfs}
\usepackage{stmaryrd}
% Les diagrammes commutatifs sont partout dans ce fonds : c'est la langue
% même de la géométrie algébrique de Grothendieck, et une transcription qui
% les rendrait en prose ne transcrirait rien.
\usepackage{tikz-cd}
% Français par défaut — c'est la langue des feuillets — mais l'anglais est
% chargé aussi : la traduction et le résumé sont des documents anglais, et la
% typographie française (espaces avant les deux-points) y serait une faute.
% Ces documents basculent par \selectlanguage{english} en tête de corps.
\usepackage[english,french]{babel}
\usepackage{fontspec}
\usepackage{geometry}
\geometry{a4paper,margin=25mm,marginparwidth=28mm}
\usepackage{marginnote}
\usepackage{xcolor}
% ulem, not soul. `soul`'s \st and \ul are built on a character-by-character
% scan that XeLaTeX's font machinery breaks: the document fails with an
% undefined control sequence rather than degrading. ulem does the same two jobs
% and is Unicode-engine safe. `normalem` stops it hijacking \emph, which the
% transcriptions use for Grothendieck's own emphasis.
\usepackage[normalem]{ulem}
\usepackage{hyperref}
\hypersetup{colorlinks=true,linkcolor=black,urlcolor=black,pdfborder={0 0 0}}

% --- Métadonnées du lot -----------------------------------------------------
% Reprises telles quelles par le script de rendu, qui les affiche en tête de
% la vue de lecture. Elles fixent ce que le fichier prétend couvrir : sans
% elles, une transcription flotte sans point d'attache dans un fonds de seize
% mille feuillets.

\newcommand{\folder}[1]{\gdef\grothFolder{#1}}
\newcommand{\batch}[1]{\gdef\grothBatch{#1}}
\newcommand{\leaves}[2]{\gdef\grothFirst{#1}\gdef\grothLast{#2}}
\newcommand{\foldertitle}[1]{\gdef\grothFoldertitle{#1}}
\gdef\grothFolder{?}\gdef\grothBatch{?}\gdef\grothFirst{?}\gdef\grothLast{?}\gdef\grothFoldertitle{}

% --- Le résumé --------------------------------------------------------------
%
% Il ouvre la lecture modernisée, sous le titre « Résumé », et il est composé
% en petit. Ce n'est pas un ornement : c'est le seul passage du document qui
% ne soit pas des mathématiques, et un lecteur doit pouvoir le reconnaître
% avant de l'avoir lu — puis le sauter s'il connaît déjà le sujet, ou s'y
% arrêter s'il ne le connaît pas. Le filet qui le suit marque la reprise du
% corps.

\newenvironment{resume}
  {\small\setlength{\parskip}{0.45em}}
  {\par\medskip\hrule\medskip}

% --- Mots-clés ---------------------------------------------------------------
%
% En anglais, en clôture du résumé de la lecture modernisée : le vocabulaire
% moderne sous lequel chercher ce que les feuillets construisent. Le manifeste
% du site (`npm run manifest`) les extrait pour étiqueter la cote entière —
% cette ligne est la source unique des tags, il n'y a pas de fichier à côté.
\newcommand{\keywords}[1]{\par\smallskip\noindent
  {\footnotesize\textsc{Keywords} --- \textit{#1}}\par}

% --- L'appareil critique ----------------------------------------------------

% Le numéro de feuillet, en marge. C'est l'ancre qui permet de régler un
% désaccord sur une formule en regardant *un* feuillet plutôt qu'une chemise
% de six cents. Le site s'en sert aussi pour tourner les pages du fac-similé
% quand on fait défiler la transcription.
\newcommand{\leaf}[1]{%
  \marginnote{\footnotesize\color{gray!70}\textsf{#1}}%
  \label{leaf:#1}\ignorespaces}

% Illisible : on ne devine pas. Un mot inventé qui se lit comme les autres est
% le pire résultat possible.
\newcommand{\ill}{\textcolor{red!60!black}{[\,\ldots\,]}}

% Lecture douteuse : le mot est donné, et signalé comme incertain.
\newcommand{\uncertain}[1]{\uline{#1}}

% Ajout de l'éditeur — une lettre manquante, un mot sous-entendu.
\newcommand{\add}[1]{\textcolor{blue!50!black}{[#1]}}

% Biffé par Grothendieck lui-même. Ce qu'il a rayé fait partie du document :
% une rature dit souvent où la pensée a changé de direction.
\newcommand{\struck}[1]{\sout{#1}}

% Note du transcripteur, sur l'état matériel du feuillet.
\newcommand{\note}[1]{\par\smallskip{\small\color{gray!60!black}\textit{[#1]}}\par\smallskip}

% Note marginale de Grothendieck, distincte du corps du texte.
\newcommand{\marginal}[1]{\par\smallskip\noindent\hspace{1em}%
  {\small\color{gray!60!black}#1}\par\smallskip}

% --- Titre ------------------------------------------------------------------

\AtBeginDocument{%
  \begin{center}
    {\large\bfseries Fonds Alexandre Grothendieck --- cote n\textsuperscript{o}~\grothFolder}\\[2pt]
    {\itshape\grothFoldertitle}\\[4pt]
    {\small feuillets \grothFirst--\grothLast\ (lot \grothBatch)}\\[2pt]
    {\footnotesize\color{gray!60!black}Transcription établie d'après le fac-similé numérisé
     par l'Université de Montpellier.\\
     Les lectures incertaines sont soulignées, les illisibles notées [\,\ldots\,],
     les ajouts entre crochets.}
  \end{center}
  \vspace{1em}}

\endinput


\folder{8}
\batch{3}
\pages{41}{60}
\dating{[vers 1967- vers 1971]}
\watermark{Édition de démonstration}
\foldertitle{Cristaux (1970). Gribouillis et calculs divers : notes manuscrites (s.d.)}

\begin{document}

\page{41}

\note{suite directe de la page 39 (lot 2) : même notation, $\varpi$,
$\Omega$, $\nu$, $r_\nu$}

En particulier, si $G_0$, $G'_0$ satisfont aux conditions de la
proposition, on a $\operatorname{Hom}(G_0, G'_0) \neq 0$, i.e.\ $G_0$ et
$G'_0$ sont isogènes, ssi $\pi = \pi'$. \struck{\ill{}} S'il en est ainsi,
alors \struck{\ill{}}
\[
\begin{cases}
\text{si } \nu \mathrel{?} \nu' : & \operatorname{Hom}_A(G_0, G'_0) \xrightarrow{\ \sim\ } \operatorname{Hom}(\Omega', \Omega) \\
\text{si } \nu > \nu' : & \operatorname{Hom}_A(G_0, G'_0) \simeq \pi \operatorname{Hom}(\Omega', \Omega)
\end{cases}
\]
\note{dans la première ligne, le signe entre $\nu$ et $\nu'$ est couvert
d'une tache d'encre ($\le$ attendu, illisible) ; après lui, « $G_0 = G'_0$ »
biffé}

\emph{N.B.} Revenons au cas général où on \uncertain{ne} suppose
\uncertain{plus} $\dim G_0 = 1$, \ill{} $\varpi$ \ill{} hauteur $= h$.
\note{le mot doublement souligné devant $\varpi$ est illisible}
Les isogénies à isogénie près \struck{\ill{}} de $G_0$, dans les $M_\nu$
($0 \le \nu \le f-1$) deviennent des vectoriels \struck{de dim.}
\uncertain{sur $K$}, isomorphes par les $\varphi_\nu$ ($0 \le \nu \le f-2$).
Alors \struck{\ill{}}
\[
M \simeq \Omega \otimes_{\mathbb{Q}_p} W = \Omega \otimes_K (K \otimes_{\mathbb{Q}_p} K_0)
\]
\note{sous $K_0$, relié par un trait : « corps des fractions de $W$ » ; les
indices $\mathbb{Q}_p$ et $K$ sont écrits sur des lettres surchargées}
et $F$ est déterminé par la connaissance de $\varphi_{f-1}$, \ill{} un
\ill{} de $M$, soit $\varpi$, \ill{} \ill{} \ill{}
\[
F_M^{\,f} = \varpi \otimes_{\mathbb{Q}_p} \mathrm{id}_{K_0} = \varpi \otimes_K (K \otimes_{\mathbb{Q}_p} K_0)
\]
La condition « d'effectivité » \uncertain{se traduit} en termes de
$\varpi$ \ill{}
\[
\begin{cases}
\varpi \text{ laisse invariant un réseau} \\
p^{f} \varpi^{-1} \text{ laisse invariant un réseau}
\end{cases}
\]
\note{l'exposant de $p$ se lit $f$ ; lecture douteuse}

\page{43}

Revenons aux conditions de la \emph{proposition}, et déterminons
\[
\omega_{G_0} = \check{t}_{G_0} \simeq M_\nu / \pi^{r} M_\nu \simeq \Omega / \pi \Omega
\]
\note{l'exposant de $\pi$ (« $r$ » ou « $F$ ») et la lettre devant le
dernier $\Omega$ ($\pi$ ou $p$) sont des lectures douteuses}
avec $k = A/\pi A$ y opérant de la façon évidente, et la structure de
vectoriel sur ($W/pW =$) $k = \mathbb{F}_q$ est déduite de la structure
précédente grâce à $\sigma^\nu$. Le cas le plus simple est celui où
$\nu = 0$, auquel cas \struck{\ill{}} les deux opérations de $k$ sur
$\omega_{G_0}$ coïncident. \struck{\ill{}} \ill{}

\note{trois lignes surchargées, en partie biffées et reliées par des traits,
se lisent : « Cas \ill{} \ill{} dictée grâce à $\sigma^\nu$ \ill{} de
$k = \mathbb{F}_q$ \struck{\ill{}}) » ; « La filtration canonique
(\ill{} $\nu \ne 1$) \ill{}
$M \otimes_W k = \sum M_\nu \otimes_A k$ \ill{} $\omega_0$ tel que
$\omega_0^{(p)} \subset M \otimes_W k^{(p)}$ soit
$M_{f-1} \otimes_A k$ \ill{} ce qui donne $\omega_0 = M_0 \otimes_A k$ » ;
« Le cas \ill{} \ill{} est »}

\[
G_0 = G_0(A, \pi) \quad \text{correspondant à } (\nu = 0 \text{ et})\ \Omega = A
\]
\ill{} \ill{} le cas d'un $\Omega$ quelconque \uncertain{par}
\[
\Omega \otimes_A G_0(A)
\]
Donc \ill{} dans ce cas comme module de Dieudonné
\[
\begin{cases}
M = A \otimes_{\mathbb{Z}_p} W \\
F_M = \mathrm{id}_A \otimes_{\mathbb{Z}_p} \sigma + \mathrm{pr}_0\bigl((\pi - 1) \otimes \sigma\bigr)
\end{cases}
\]
où $\mathrm{pr}_0$ est le projecteur sur la composante correspondant au
$A$-hom
\[
A \otimes_{\mathbb{Z}_p} W \to A
\]
qui sur $W$ est l'inclusion $W \hookrightarrow A$.

\note{un double trait sépare la suite}

Nous nous proposons maintenant de relever $G_0$ \struck{l'action} en un
groupe de BT $G$ sur $A$, avec action de $A$ sur $G$, de telle façon que
l'action \uncertain{correspondante} de $A$ sur $t_G$ (qui est un $A$-module
inversible) soit l'action évidente. D'après Lubin--Tate,

\page{45}

un tel relèvement existe et est unique (à isom.\ unique près).
\struck{\ill{}} Pour simplifier et pouvoir appliquer la théorie des
déformations, on \uncertain{suppose} que $\pi A$ est stable par puissances
divisées. \ill{} relever $G_0$ en un groupe de B.T.\ $G$ sur $A$
(\uncertain{à isom.\ près}, \ill{} \ill{} de l'action de $A$) revient à
relever la filtration de
\[
M \otimes_W k = M_A \otimes_A k \ \text{ en une filtration de } M_A \ (\text{sur } A).
\]
\note{« (à isom.\ unique près) » est ajouté au-dessus de la ligne ; la
lecture de la parenthèse de la cinquième ligne est douteuse}

\note{encadré à gauche : « $\Omega$ \quad
$M = \Omega \otimes_{\mathbb{Z}_p} W$, \quad
$M_A = \Omega \otimes_{\mathbb{Z}_p} A$ » et le diagramme}

\begin{tikzcd}[column sep=small, row sep=small]
A \arrow[r, no head] & A \otimes_{\mathbb{Z}_p} W \arrow[r, no head] & A \otimes_{\mathbb{Z}_p} A \\
\mathbb{Z}_p \arrow[u, no head] \arrow[r, no head] & W \arrow[u, no head] \arrow[r, no head] & A \arrow[u, no head]
\end{tikzcd}

La compatibilité à l'action de $A$ signifie qu'il s'agit d'une filtration
de $A \otimes_{\mathbb{Z}_p} A$-modules, i.e.\ d'un sous-$A \otimes_{\mathbb{Z}_p} A$-module
\uncertain{de rang 1 sur $1 \otimes A$} \struck{\ill{}} de
$\Omega \otimes_{\mathbb{Z}_p} A$, (le dernier \ill{} \struck{\ill{}})
d'ailleurs can.\ \uncertain{isom.}\ à $\omega_G$. On a
\[
0 \to \omega_G \to \Omega \otimes_{\mathbb{Z}_p} A \to t_{G^*} \to 0
\]
et, par transposition
\[
(**) \qquad 0 \to \omega_{G^*} \to \check{\Omega} \otimes_{\mathbb{Z}_p} A \xrightarrow{\ p_0\ } t_G \to 0
\qquad\qquad A \otimes_{\mathbb{Z}_p} A \xrightarrow{\ p_0\ } A
\]
\note{au-dessus de la flèche vers $t_G$, une lettre surchargée par
« $p_0$ »}

L'hypothèse que \uncertain{sur $\omega_G$} l'action de $A$ (via l'action
sur $G$) est l'action de $A$-module naturelle, signifie que \struck{\ill{}}
$\omega_G$ (\uncertain{ou} $t_G$) \ill{} \ill{} \struck{et
$A \otimes_{\mathbb{Z}_p} A$ \ill{} \ill{} $p_0$ \ill{}}, \ill{} \ill{} de
$A$-module \struck{trivial} évident par \ill{} \uncertain{particulier} les
relations : \uncertain{à l'aide de}
\[
p_0 : A \otimes_{\mathbb{Z}_p} A \to A, \qquad p_0(\lambda \otimes \mu) = \lambda\mu
\]

\page{47}

Cela, plus le fait que $t_G$ et $\check{\Omega} \otimes_{\mathbb{Z}_p} A$ sont
de rang 1 sur $A$ via $A \otimes_{\mathbb{Z}_p} A$, et que $p$ est surjectif,
prouve que $t_G$ est \emph{déduit de} $\check{\Omega} \otimes_{\mathbb{Z}_p} A$
par le changement de base $A \otimes_{\mathbb{Z}_p} A \xrightarrow{p_0} A$,
\struck{\ill{}} (comme $\check{\Omega} \otimes_{\mathbb{Z}_p} A$ est déduit de
$\check{\Omega}$ par $A \xrightarrow{\lambda \mapsto \lambda \otimes 1}
A \otimes_{\mathbb{Z}_p} A$) --- d'où un isom.\ can.
\[
t_G \simeq \check{\Omega} \qquad \text{i.e.} \qquad \omega_G \simeq \Omega
\]
et la suite exacte (**) et la suite exacte déduite (*) par dualité sont
canoniques évidentes, (\ill{} \ill{} cette filtration \ill{} \ill{}
vectorielle \ill{}) étend celle de $M_A \otimes_A k$, \ill{} \ill{} le th.\
d'existence et d'unicité de Lubin--Tate, \struck{\ill{}} sous une forme moins
forte d'ailleurs en ce que \ill{} l'unicité, et sous l'hyp.\
complémentaire que $\exists$ puiss.\ div.\ sur $\pi A$.
\note{« dit (*) par dualité » est ajouté au-dessus de la ligne et relié
par un trait}

Si on abandonne cette dernière hypothèse, et qu'on utilise la
classification des déformations à isogénie près, on trouve de
\uncertain{même} \ill{} une suite exacte canonique
\[
0 \to \omega_G \otimes_A K \to M_K \to t_{G^*} \otimes_A K \to 0,
\qquad M_K = \Omega \otimes_{\mathbb{Z}_p} K
\]
\ill{} \uncertain{aussi}
\[
0 \to \omega_{G^*} \otimes_A K \to \check{M}_K \to t_G \otimes_A K \to 0,
\qquad \check{M}_K = \check{\Omega} \otimes_{\mathbb{Z}_p} K
\]
\note{les égalités $M_K = \ldots$, $\check{M}_K = \ldots$ sont écrites
verticalement sous $M_K$, $\check{M}_K$ ; sous $M_K$, un $K$ surchargé}
et on trouve \ill{} que
\[
t_G \otimes_A K \simeq \check{M}_K \otimes_{\check{\Omega} \otimes_{\mathbb{Z}_p} K} (K \otimes K, p_0) \simeq \check{\Omega} \otimes_A K
\]
\note{l'indice du produit tensoriel du milieu est une lecture douteuse}

\page{49}

ce qui redonne la filtration de $\check{M}_K$ (ou de façon équivalente, de
$M_K$) en termes de $\Omega$. Cela donne bien th.\ d'unicité à isogénie
près (\uncertain{mais} pas d'existence ni d'unicité à \emph{isom.} près,
pour lesquelles il faut ici se reporter à Lubin--Tate). Notons que $t_G$ et
$\check{\Omega}$ s'identifient à un sous-réseau de rang 1 \ill{} d'un
vectoriel de dim 1 sur $K$, donc (utilisant le choix de $\pi \in A$) ils
sont can.\ isomorphes, \struck{\ill{}} par $x \mapsto \pi^{-\alpha} x$
($\alpha \in \mathbb{Z}$ \ill{}),
\[
t_G = \pi^{-\alpha} \check{\Omega} \qquad \text{i.e.} \qquad \Omega = \pi^{\alpha} \omega_G,
\]
\struck{On peut} où $\alpha \in \mathbb{Z}$ est bien déterminé (par $A$ et
$\pi$, $\alpha = \ldots$ \ill{} $A$ \uncertain{tous} seuls comme \ill{} \ill{}
\ldots).

On a aussi
\[
\mathbb{D}^*(G)_K \simeq M_K \simeq \Omega \otimes_{\mathbb{Z}_p} K
\]
et \uncertain{notons} que $\mathbb{D}^*(G)$ et $M$ s'identifient à des
sous-réseaux d'un même vectoriel sur $K$. On peut se demander si ces
sous-réseaux sont les mêmes (comme dans le cas où $\exists$ puiss.\ div.\
sur $\pi A$), ce qui impliquerait évidemment $\alpha = 0$.. Cela ne dépend
pas du choix de $\pi$, mais seulement de $A$ \ldots\ldots

\page{51}

Nous allons examiner \struck{\ill{}} ce que donne le théorème de Tate, pour
le groupe de Lubin--Tate $G$. Notons que, si $K_\pi$ ($\subset \bar{K}$) est
l'extension \emph{abélienne} de groupe $A^*$ définie par $G$ (via
$\operatorname{Gal}(\bar{K}/K) \xrightarrow{\chi_\pi} A^*$), alors
\struck{\ill{}} $T_p(G)(\bar{K})$ est un module inversible sur $A$, et les
bases de celui-ci se permutent de façon simplement transitive par
$\operatorname{Gal}(K_\pi/K)$. Donc, il y a un choix privilégié (à isom.\
\emph{unique} près) d'une extension algébrique $K_\pi$ de $K$ (indépendant
du choix d'une clôture $\bar{K}$ !) \uncertain{avec} un générateur
\[
e \in T_p(G)^*(K_\pi)
\]
[en fait, $\operatorname{Spec} K_\pi$ représente le \uncertain{faisceau}
$T_p(G)^* = \varprojlim G(\nu)^*$, où la $*$ désigne les générateurs
\struck{\ill{}} $A$-modules \ldots]. Nous supposons fixé $K_\pi$
\uncertain{de cette} \uncertain{façon}, et choisissons pour $\bar{K}$ une
\uncertain{sur-extension} de $K_\pi$. Donc le choix \struck{\ill{}} ci-dessus
donne un isom.\ de $A$-modules
\[
\overbrace{T_p(G)(K_\pi)}^{T_p} \xrightarrow{\ \sim\ } A, \qquad \lambda \mapsto \lambda e .
\]

\note{encadré à gauche :}

\begin{tikzcd}[column sep=small, row sep=small]
T_p \simeq A \arrow[r, no head] & T_p \otimes A \simeq A \otimes_{\mathbb{Z}_p} A \arrow[r, no head] & T_p \otimes C \simeq A \otimes_{\mathbb{Z}_p} C \\
\mathbb{Z}_p \arrow[u, no head] \arrow[r, no head] & A \arrow[u, no head] \arrow[r, no head] & C \arrow[u, no head]
\end{tikzcd}

On regarde sur $T_p(\ldots) \otimes_A C$ d'une part l'action de $A$ provenant
de l'action de $A$ sur $G$ (« action via opérateurs »), \uncertain{d'autre}
part l'action provenant de la structure de $C$-vectoriel. Enfin, on a
l'action de $\boldsymbol{\pi} = \operatorname{Gal}(\bar{K}/K)$, qui commute à
l'action (\uncertain{les deux}) de $A$. Supposons $K/\mathbb{Q}_p$
galoisienne pour simplifier.
\note{dans $T_p(\ldots)$ une lettre est surchargée ; le $\boldsymbol{\pi}$
du groupe de Galois est écrit en gras, distinct de l'uniformisante $\pi$ ;
on le garde en gras sur les pages 51 à 55}

\page{53}

Alors comme module sur $A_C$, \struck{$T_p \otimes C$ s'identifie}
canoniquement
\[
(*) \qquad \boxed{\ T_p \otimes_{\mathbb{Z}_p} C \simeq \prod_\sigma C_\sigma\ }
\]
où $C_\sigma = C$, comme $C$-vectoriel, l'opération de $A$ donnée par la
\uncertain{façon} via $\sigma : A \to C$, où
\[
* \qquad \sigma \in \operatorname{Hom}_{\mathbb{Q}_p}(K, C) \simeq \operatorname{Aut}_{\mathbb{Q}_p}(K) \simeq \operatorname{Aut}_{\mathbb{Z}_p}(A) \simeq \operatorname{Hom}_{\mathbb{Z}_p}(A, C).
\]
\note{le premier « Hom » est écrit sur un mot surchargé ; au-dessus de la
première ligne biffée, « \ill{} \ill{} \ill{} »}

Cette décomposition est compatible avec l'action de $\boldsymbol{\pi}$, de
façon précise, $\boldsymbol{\pi}$ opère sur le facteur $C_\sigma$ par
$\rho_\sigma(g)$ ($g \in \boldsymbol{\pi}$) donnés par
\[
x \in C_\sigma \Rightarrow \rho_\sigma(g) x = \sigma(\chi_\pi(g)) x
\]
Comparons cette décomposition avec celle donnée par le théorème de Tate

\begin{tikzcd}[column sep=small]
0 \arrow[r] & \mathfrak{t} \otimes C(1) \arrow[r] & T_p \otimes C \arrow[r] \arrow[l, dashed, bend left=30] & \check{\mathfrak{t}}^{*} \otimes C \arrow[r] \arrow[l, dashed, bend left=30] & 0
\end{tikzcd}

\[
(**) \qquad T_p \otimes C \simeq \mathfrak{t} \otimes_A C(1) + \check{\mathfrak{t}}^{*} \otimes_A C
\]
\note{la lettre notée $\mathfrak{t}$ (ici et page 55) est une majuscule
cursive à longue hampe, de lecture douteuse ; les deux flèches en pointillé
sont des arcs sous la suite, une accolade réunit la suite et (**)}

Notons que ce dernier isom.\ est \uncertain{fonctoriel}, donc commute à
l'action de $A$ sur les deux membres. Or l'action de $A$ sur $\mathfrak{t}$
est l'action régulière, \uncertain{i.e.} \ill{} \ill{} que, et l'action de
$A$ sur $\check{\mathfrak{t}}^{*} \otimes C$ se fait suivant les hom
$A \to C$ pour $\sigma \neq 1$. Donc la décomposition (**) est celle
déduite de la décomposition de $T_p \otimes C$ \struck{\ill{} \ill{}} en
l'écrivant
\[
T_p \otimes C \simeq C_1 \oplus \prod_{\sigma \neq 1} C_\sigma .
\]
\struck{Il faut} En tenant \uncertain{compte} \uncertain{des} isomorphismes
can., et

\page{55}

tenant compte des isom.\ can.\ obtenus plus haut (de $A \otimes C$-modules)
\[
\mathfrak{t} \otimes_A C \simeq C_1, \qquad
\check{\mathfrak{t}}^{*} \otimes_A C \simeq \prod_{\sigma \neq 1} C_\sigma,
\]
\struck{il faut donc définir} le théorème de Tate équivaut à la donnée des
isom.\ de $A \otimes_{\mathbb{Z}_p} C$-modules
\[
C_\sigma \to C_\sigma, \quad \text{i.e.} \quad \lambda \mapsto \lambda x_\sigma \quad (x_\sigma \in C)
\]
qui commutent à l'action de $\boldsymbol{\pi}$ opérant sur la source par
$*$, et sur le but par $x \mapsto \tau(x)\, g(x)$ si $\sigma = 1$, par
$x \mapsto g(x)$ si $\sigma \neq 1$ ($\tau$ étant le caractère de Tate).
Cela revient aux conditions,
\[
\boxed{\ \begin{cases}
g(x_\sigma) = \sigma(\chi_\pi(g))\, x_\sigma & \text{si } \sigma \neq 1 \\
g(x_1) = \tau(g)^{-1} \chi_\pi(g)\, x_1 & \text{si } \sigma = 1
\end{cases}
\quad \text{si } g \in \boldsymbol{\pi} = \operatorname{Gal}(\bar{K}/K)\ }
\]
\note{« = $\operatorname{Gal}(\bar K/K)$ » est écrit sous $\boldsymbol{\pi}$
; l'encadré est ouvert à gauche}

\emph{NB} A priori, les $x_\sigma$ sont déterminés par la condition à une
multiplication près par un élément de certains modules $A^* \subset K^*
\subset C^*$. Mais la \ill{} \uncertain{restriction} \ill{} \ill{} \ill{}
$\bar{K}$) le choix de l'un des $x_\sigma$ (disons, de $x_1$)
\uncertain{détermine} canoniquement, par le théorème de Tate, tous les
autres. \struck{\ill{}} Il n'existe pas de $x'_1 \in C^*$ tel qu'on ait
\[
g(x'_1) = \chi_\pi(g)\, x'_1 \qquad \text{pour } g \in \boldsymbol{\pi}
\ (\text{ou} = \text{pour } g \in \text{Inertie de } \boldsymbol{\pi})
\]
\uncertain{car} en faisant le produit de $x'_1$ et des $x_\sigma$
($\sigma \neq 1$), on trouverait $x \in C^*$ tel que
\[
g(x) = \underbrace{N_{A^*/\mathbb{Z}_p^*}(\chi_\pi(g))}_{= \tau(g) \text{ si } g \in \text{Inertie}}\, x ,
\]
impossible.
\note{la ligne « Mais la \ldots\ $\bar K$) le choix de » est très cursive ;
seuls quelques mots se lisent}

\page{57}

\note{feuillet de papier quadrillé tenu en largeur, au crayon. En haut à
gauche, sous l'en-tête « $n/17$ », une colonne de petites lignes polygonales,
une par ligne : « $n = 1, 2, 3$ », « $n = 4$ », « $n = 5$ », « $n = 6$ »,
« $n = 7, 8, 9, 10, 11, 12, 13$ », « $n = 14$ », « $n = 15, 16$ » (le 16
surchargé) ; à côté, une colonne titrée « $\nu$ » porte $1, 2, 3, 1, 2, 3,
2$. Au centre, deux longues droites parallèles partent de l'origine, avec
des points marqués ; à droite, un faisceau de segments aboutit à un axe
vertical gradué de 1 à 16 ; l'axe horizontal est gradué de 1 à 18}

Inscriptions, de haut en bas et de gauche à droite :
\[
\begin{gathered}
\ldots\ \tfrac{6}{9} > \tfrac{11}{17}, \qquad \tfrac{2}{3} > \tfrac{11}{17}, \qquad
\tfrac{5}{12} > \tfrac{7}{17}, \qquad \tfrac{1}{4} < \tfrac{?}{11} < \tfrac{2}{7}, \qquad
\tfrac{2}{7} < \tfrac{5}{17}, \\
\tfrac{4}{17} < \tfrac{1}{4}, \qquad \tfrac{1}{9} < \tfrac{2}{17}, \qquad
\tfrac{1}{6} < \tfrac{3}{17}, \qquad \tfrac{7}{16} < \tfrac{12}{17}, \qquad
\tfrac{5}{7} \mathrel{?} \tfrac{13}{18}, \\
\tfrac{4}{13} > \tfrac{7}{23}, \qquad {<}\ \tfrac{51}{13}\,(?), \qquad
\tfrac{3}{10} > \tfrac{7}{24}, \qquad \tfrac{2}{7} < \tfrac{7}{24}, \qquad
\tfrac{1}{3} > \tfrac{7}{23}, \qquad \tfrac{2}{7} < \tfrac{1}{3}, \\
\tfrac{7}{9} < \tfrac{7}{8}\,(?), \qquad \tfrac{4}{5} < \tfrac{14}{17}\,(?), \qquad
\tfrac{11}{17} > \tfrac{13}{23}, \qquad \tfrac{?}{13} > \tfrac{15}{23}, \qquad
\tfrac{10}{14} < \tfrac{13}{18}, \\
\tfrac{13}{17} < \tfrac{10}{13}, \qquad \tfrac{3}{5} < \tfrac{14}{23}, \qquad
\tfrac{3}{4} < \tfrac{13}{17}
\end{gathered}
\]
\note{« $4/17 < 1/4$ » : le signe est corrigé à l'encre sur un autre ; dans
« $5/7 \mathrel{?} 13/18$ » le signe est surchargé ; « $2/7 < 7/24$ » est
précédé d'une fraction biffée ; « $2/7 < 1/3$ » est écrit sur une première
fraction ; « $11/17 > 13/23$ » et « $\ldots/13 > 15/23$ », encadrés
ensemble, avec un « ? » entre eux ; près du centre, « $2533 > 285$ »
(lecture douteuse) ; un gribouillis en haut à droite. Les inégalités sont
transcrites telles qu'écrites, sans vérification}

\page{58}

\note{même papier quadrillé, au crayon : deux droites presque parallèles
partant de l'origine, l'une passant par une suite de points marqués, sur
l'axe horizontal tracé jusqu'à l'abscisse 16 environ}

\[
\tfrac{4}{5} < \tfrac{13}{15}, \qquad \tfrac{12}{14} = \tfrac{6}{7}, \qquad
\tfrac{14}{16} = \tfrac{7}{8} <
\]

\page{59}

\note{même papier, tenu en hauteur, au crayon : en haut, deux droites
partant de deux points voisins de l'origine, l'une par une suite de points
marqués ; plus bas, un second faisceau de six droites de pentes
décroissantes, issues de deux points voisins, dont une ligne brisée au
départ}

\[
\tfrac{12}{14} = \tfrac{6}{7}, \qquad \tfrac{13}{15} < \tfrac{7}{8}, \qquad
\tfrac{3}{4} > \tfrac{11}{15}, \qquad \tfrac{2}{3} < \tfrac{11}{16}
\]
\note{une quatrième inscription, très pâle, sous les précédentes, est
illisible}

\page{60}

\marginal{$\mathfrak{G}$}\note{en marge gauche, un $\mathfrak{G}$ gothique
repassé à l'encre}

Groupe de Galois d'une catégorie tannakienne
\uncertain{(gr.\ déri.)} ($\mathbb{Q}_p$-linéaire, ½ simple) $\mathscr{C}$
munie d'un foncteur fibre $\Delta$ : \ill{}
dans les $F$-cristaux (filtrés \uncertain{gr.\ déri.}) à isogénie près sur
un anneau \struck{local} de val.\ discr.\ complet $V$ d'inégales car., à corps
résiduel parfait $k$, corps des fractions $L$.
\note{« gr.\ déri.\ » (deux fois), « ($\mathbb{Q}_p$-linéaire) »,
« ½ simple » (i.e.\ semi-simple), « $F$- » et « filtrés » sont des ajouts
interlinéaires ; $\Delta$ est écrit au-dessus de « fibre »}

On a

\begin{tikzcd}[column sep=small, row sep=small]
W \arrow[r, hook] \arrow[d, hook] & V \arrow[d, hook] \\
K \arrow[r, hook] & L
\end{tikzcd}

\note{« vect.\ Witt » écrit contre $W$}

\begin{enumerate}
\item[a)] foncteur \uncertain{fibre $\Phi$} vers les $F$-cristaux sur $k$, à
isogénie près $=$ certains couples $(E, F)$, où $E$ vectoriel de dim finie
sur $K$, $F$ endomorphisme $\sigma$-linéaire de $E$, tel qu'il existe un
réseau invariant par $F$.
\item[b)] filtration de $\Phi \otimes_K L$.
\end{enumerate}

\struck{Le foncteur $\Phi$ (permet d'identifier $\mathscr{C}$ à la catégorie
des repr.\ \ill{} sur $K$ \ill{}) (en utilisant la \ill{})}\note{passage
encadré et biffé de hachures obliques}

Le foncteur fondamental fournit un foncteur
\[
T_p : \mathscr{C} \longrightarrow \operatorname{Mod}_f(\pi, \mathbb{Q}_p)
\qquad \text{où } \pi = \operatorname{Gal}(\bar{L}/L)
\]
donc \emph{un foncteur fibre sur} $\mathbb{Q}_p$, correspondant à une
représentation (à isom.\ près) de $\mathscr{C}$ comme
\[
\mathscr{C} \cong \operatorname{Rep}_{\mathbb{Q}_p}(G)
\]
$G$ étant un groupe pro-\struck{\ill{}} \ill{} sur $\mathbb{Q}_p$.
\note{« $G$ » est encadré ; le mot après « pro- » est biffé et récrit, les
deux illisibles}

\end{document}
